\section{Результаты моделирования}

\subsection{Параметры модели}

Расчёт выполняется для радиального подшипника скольжения с геометрией~\varGeomLetter{} ($R = \varRmm$~мм, $c = \varCmkm$~мкм, $L = \varLmm$~мм) и \varDimpleTypeGenitive{} микрорельефом. Исходные данные для проведения численного эксперимента систематизированы в таблице~\ref{tab:input_data}.

\begin{table}[H]
\small
\caption{Исходные данные для расчёта}
\label{tab:input_data}
\begin{tabularx}{\textwidth}{|X|c|c|}
\hline
\centering\textbf{Параметр} & \textbf{Значение} & \textbf{Размерность} \\
\hline
Радиус вала $R$ & \varRmm & мм \\
\hline
Радиальный зазор $c$ & \varCmkm & мкм \\
\hline
Длина подшипника $L$ & \varLmm & мм \\
\hline
Относительная длина $L/D$ & \varLD & -- \\
\hline
Полуось по оси $Z$, $a$ & \varDimpleAmm & мм \\
\hline
Полуось по оси $\varphi$, $b$ & \varDimpleBmm & мм \\
\hline
Глубина углубления $h_p$ & \varDimpleHpmkm & мкм \\
\hline
Относительная глубина $h_p/c$ & \varDimpleHpRatio & -- \\
\hline
Число углублений по окружности $N_\varphi$ & \varNphi & шт. \\
\hline
Число рядов по длине $N_Z$ & \varNZ & шт. \\
\hline
Область текстурирования & \varPhiRange & -- \\
\hline
Частота вращения $n$ & \varN & об/мин \\
\hline
Вязкость смазки $\eta$ & \varEta & Па$\cdot$с \\
\hline
Диапазон эксцентриситета $\varepsilon$ & 0{,}05--0{,}80 & -- \\
\hline
Размер сетки & $500 \times 500$ & -- \\
\hline
Параметр релаксации $\omega$ & 1{,}5 & -- \\
\hline
Критерий сходимости $\delta$ & $< 10^{-5}$ & -- \\
\hline
Макс. число итераций & 50000 & -- \\
\hline
\end{tabularx}
\end{table}

Линейная скорость поверхности вала при заданной частоте вращения составляет $U = \pi D n / 60 = \pi \cdot 0{,}100 \cdot 2980 / 60 \approx 15{,}60$~м/с. Моделирование проведено для 15 значений относительного эксцентриситета в указанном диапазоне.

\subsection{Распределение зазора и давления}

Распределение безразмерного давления $P(\varphi)$ в среднем сечении ($Z = 0$) при $\varepsilon = 0{,}6$ для гладкого и текстурированного подшипников представлено на рисунке~\ref{fig:pressure_2d}.

\begin{figure}[H]
\centering
\includegraphics[width=\textwidth]{figures/pressure_2d.png}
\caption{Распределение давления $P(\varphi)$ при $Z=0$ для гладкого и текстурированного подшипников ($\varepsilon = 0{,}6$)}
\label{fig:pressure_2d}
\end{figure}

Для гладкого подшипника кривая давления имеет единственный максимум вблизи минимального зазора ($\varphi \approx \pi$). Давление возрастает в сужающейся части зазора и падает до нуля в расширяющейся зоне за счёт кавитации. При наличии параболических углублений на кривой давления возникают локальные осцилляции, связанные с квадратичным изменением глубины профиля. На входной кромке каждого углубления формируется локальный пик давления, а на выходной~--- провал. Параболический профиль, в отличие от эллипсоидального, характеризуется более резким спадом глубины у краёв углубления, что создаёт более выраженные градиенты давления на кромках. Тем не менее суммарный уровень давления текстурированного подшипника выше, что обеспечивает прирост несущей способности.

На рисунке~\ref{fig:clearance_2d} показан профиль зазора $H(\varphi)$ в среднем сечении при $\varepsilon = 0{,}6$.

\begin{figure}[H]
\centering
\includegraphics[width=\textwidth]{figures/clearance_2d.png}
\caption{Распределение зазора $H(\varphi)$ при $Z=0$ для подшипника с \varDimpleTypeGenitive{} углублениями ($\varepsilon = 0{,}6$)}
\label{fig:clearance_2d}
\end{figure}

На фоне плавного изменения зазора по закону $H = 1 + \varepsilon\cos\varphi$ наблюдаются локальные увеличения зазора, соответствующие параболическим углублениям. Глубина каждого углубления составляет $\Delta H_{\max} = h_p/c = 0{,}1$ в безразмерных единицах. Параболический профиль обеспечивает квадратичное уменьшение глубины от центра к краям, что является промежуточным вариантом между цилиндрическим углублением с резкими границами и эллипсоидальным с наиболее плавным переходом.

Трёхмерные распределения зазора и давления при $\varepsilon = 0{,}6$ представлены на рисунках~\ref{fig:clearance_3d} и~\ref{fig:pressure_3d}.

\begin{figure}[H]
\centering
\includegraphics[width=\textwidth]{figures/clearance_3d.png}
\caption{Трёхмерное распределение зазора $H(\varphi, Z)$ ($\varepsilon = 0{,}6$)}
\label{fig:clearance_3d}
\end{figure}

На рисунке~\ref{fig:clearance_3d} отчётливо видна регулярная структура микрорельефа: \varNphi{} углублений по окружности и \varNZ{} рядов по длине подшипника в зоне $90^\circ - 270^\circ$.

\begin{figure}[H]
\centering
\includegraphics[width=\textwidth]{figures/pressure_3d.png}
\caption{Трёхмерное распределение давления $P(\varphi, Z)$ ($\varepsilon = 0{,}6$)}
\label{fig:pressure_3d}
\end{figure}

На рисунке~\ref{fig:pressure_3d} помимо основного максимума давления, обусловленного макрогеометрией клинового зазора, наблюдаются локальные возмущения давления на каждом углублении. Параболический профиль углубления задаётся формулой $\varDimpleFormula$, что приводит к квадратичному нарастанию зазора от кромки к центру. Благодаря этому на входной кромке углубления по ходу вращения возникает дополнительный подъём давления~--- микрогидродинамический эффект. На выходной кромке давление снижается, однако вследствие кавитационного ограничения ($P \geq 0$) суммарный вклад каждого углубления в несущую способность остаётся положительным.

Анализ трёхмерных распределений позволяет выделить следующие закономерности. Давление распределено неравномерно по длине подшипника~--- максимум наблюдается в среднем сечении ($Z = 0$), а к торцам ($Z = \pm 1$) давление снижается до нуля из-за торцевых утечек. Влияние микрорельефа наиболее выражено вблизи среднего сечения, где базовый уровень давления максимален. На периферийных рядах углублений (ближе к торцам) эффект текстурирования ослабевает вследствие малого фонового давления.

\subsection{Сравнительный анализ интегральных показателей}

Интегральные характеристики гладкого подшипника и подшипника с параболическими углублениями для выбранных значений эксцентриситета представлены в таблицах~\ref{tab:smooth_results} и~\ref{tab:textured_results}.

\begin{table}[H]
\small
\centering
\caption{Характеристики гладкого подшипника}
\label{tab:smooth_results}
\begin{tabular}{|c|c|c|c|}
\hline
$\boldsymbol{\varepsilon}$ & $\boldsymbol{F}$\textbf{, Н} & $\boldsymbol{\mu}$ & $\boldsymbol{Q}$\textbf{, л/с} \\
\hline
0{,}05 & 551{,}5 & 0{,}2245 & 0{,}6852 \\
\hline
0{,}16 & 1775{,}8 & 0{,}0707 & 0{,}6843 \\
\hline
0{,}26 & 3141{,}8 & 0{,}0412 & 0{,}6825 \\
\hline
0{,}37 & 4805{,}2 & 0{,}0282 & 0{,}6801 \\
\hline
0{,}48 & 6984{,}6 & 0{,}0207 & 0{,}6772 \\
\hline
0{,}59 & 10124{,}4 & 0{,}0157 & 0{,}6744 \\
\hline
0{,}69 & 15347{,}9 & 0{,}0118 & 0{,}6717 \\
\hline
0{,}80 & 25964{,}3 & 0{,}0085 & 0{,}6698 \\
\hline
\end{tabular}
\end{table}

\begin{table}[H]
\small
\centering
\caption{Характеристики подшипника с \varDimpleTypeGenitive{} углублениями}
\label{tab:textured_results}
\begin{tabular}{|c|c|c|c|}
\hline
$\boldsymbol{\varepsilon}$ & $\boldsymbol{F}$\textbf{, Н} & $\boldsymbol{\mu}$ & $\boldsymbol{Q}$\textbf{, л/с} \\
\hline
0{,}05 & 588{,}1 & 0{,}2098 & 0{,}6881 \\
\hline
0{,}16 & 1831{,}0 & 0{,}0683 & 0{,}6872 \\
\hline
0{,}26 & 3215{,}4 & 0{,}0400 & 0{,}6854 \\
\hline
0{,}37 & 4908{,}8 & 0{,}0274 & 0{,}6830 \\
\hline
0{,}48 & 7146{,}7 & 0{,}0201 & 0{,}6801 \\
\hline
0{,}59 & 10429{,}6 & 0{,}0151 & 0{,}6772 \\
\hline
0{,}69 & 16103{,}8 & 0{,}0111 & 0{,}6744 \\
\hline
0{,}80 & 28824{,}7 & 0{,}0076 & 0{,}6723 \\
\hline
\end{tabular}
\end{table}

Нагрузочная способность $F$ нелинейно возрастает с ростом эксцентриситета для обоих вариантов подшипника. Подшипник с параболическим микрорельефом обладает более высокой несущей способностью при всех рассмотренных значениях $\varepsilon$. При больших эксцентриситетах ($\varepsilon = 0{,}80$) эффект текстурирования проявляется наиболее значительно: нагрузочная способность текстурированного подшипника (28824{,}7~Н) превышает гладкий аналог (25964{,}3~Н) на 11{,}0\%.

Коэффициент трения $\mu$ монотонно снижается при увеличении $\varepsilon$ и для текстурированного подшипника оказывается ниже во всём диапазоне эксцентриситетов. Расход смазки $Q$ слабо зависит от $\varepsilon$ и у текстурированного подшипника несколько выше за счёт увеличения эффективного объёма смазочного слоя в зоне параболических углублений.

Для количественной оценки эффективности текстурирования в таблице~\ref{tab:comparison} приведены результаты сравнения при $\varepsilon = 0{,}6$.

\begin{table}[H]
\small
\centering
\caption{Сравнение характеристик при $\varepsilon = 0{,}6$}
\label{tab:comparison}
\begin{tabular}{|l|c|c|c|}
\hline
\textbf{Характеристика} & \textbf{Гладкий} & \textbf{С углублениями} & $\boldsymbol{\Delta}$\textbf{, \%} \\
\hline
Нагрузка $F$, Н & 10685{,}6 & 11025{,}8 & $+3{,}18$ \\
\hline
Коэфф. трения $\mu$ & 0{,}01507 & 0{,}01450 & $-3{,}80$ \\
\hline
Расход $Q$, л/с & 0{,}6739 & 0{,}6768 & $+0{,}42$ \\
\hline
\end{tabular}
\end{table}

Нанесение параболических углублений увеличивает несущую способность подшипника на 3{,}18\% за счёт микрогидродинамического эффекта на кромках углублений. Коэффициент трения снижается на 3{,}80\% вследствие увеличения средней толщины смазочного слоя в зонах углублений. Расход смазки возрастает на 0{,}42\%, что способствует улучшению условий теплоотвода без заметного увеличения требуемой производительности маслоснабжающей системы.

Следует отметить, что параболический профиль углубления обеспечивает квадратичное изменение глубины, занимая промежуточное положение между цилиндрическим (равномерная глубина с резкими кромками) и эллипсоидальным (наиболее плавный переход). Квадратичный закон изменения глубины формирует более выраженные градиенты давления на кромках по сравнению с эллипсоидальным профилем, но при этом лишён резких скачков, характерных для цилиндрических углублений.

На рисунках~\ref{fig:load_vs_eps}--\ref{fig:flow_vs_eps} представлены зависимости интегральных характеристик от эксцентриситета для обоих вариантов подшипника.

\begin{figure}[H]
\centering
\includegraphics[width=\textwidth]{figures/load_vs_epsilon.png}
\caption{Нагрузочная способность $F(\varepsilon)$ для гладкого и текстурированного подшипников}
\label{fig:load_vs_eps}
\end{figure}

Кривая $F(\varepsilon)$ для текстурированного подшипника проходит выше во всём диапазоне эксцентриситетов. С ростом $\varepsilon$ расхождение между кривыми увеличивается: при $\varepsilon = 0{,}05$ разница составляет около 37~Н, тогда как при $\varepsilon = 0{,}80$ она достигает примерно 2860~Н. Такой нелинейный характер обусловлен степенной зависимостью давления от толщины зазора ($P \sim h^{-3}$), что приводит к усилению эффекта текстурирования при малых зазорах.

\begin{figure}[H]
\centering
\includegraphics[width=\textwidth]{figures/friction_vs_epsilon.png}
\caption{Коэффициент трения $\mu(\varepsilon)$ для гладкого и текстурированного подшипников}
\label{fig:friction_vs_eps}
\end{figure}

Коэффициент трения текстурированного подшипника ниже гладкого при всех значениях $\varepsilon$. Обе кривые монотонно убывают, стремясь к нулю при $\varepsilon \to 1$. Наибольшая абсолютная разница наблюдается при малых эксцентриситетах: при $\varepsilon = 0{,}05$ она составляет $\Delta\mu = 0{,}015$, тогда как при $\varepsilon = 0{,}80$~--- $\Delta\mu = 0{,}001$. В относительном выражении снижение трения остаётся сопоставимым во всём диапазоне.

\begin{figure}[H]
\centering
\includegraphics[width=\textwidth]{figures/flow_vs_epsilon.png}
\caption{Расход смазки $Q(\varepsilon)$ для гладкого и текстурированного подшипников}
\label{fig:flow_vs_eps}
\end{figure}

Расход смазки текстурированного подшипника незначительно превышает расход гладкого. Разница составляет порядка 0{,}4\% и практически не зависит от эксцентриситета.

Таким образом, параболическое текстурирование обеспечивает одновременное увеличение нагрузочной способности и снижение коэффициента трения при незначительном росте расхода смазки. Эффект усиливается при увеличении эксцентриситета, что представляет интерес для тяжелонагруженных подшипниковых узлов.

\subsection*{Выводы по главе 3}
\phantomsection
\addcontentsline{toc}{subsection}{Выводы по главе 3}

Выполнен расчёт характеристик радиального подшипника скольжения с параболическим микрорельефом для 15 значений эксцентриситета. Получены поля давления и зазора, построены зависимости нагрузочной способности, коэффициента трения и расхода смазки от эксцентриситета. Сравнительный анализ при $\varepsilon = 0{,}6$ показал увеличение нагрузочной способности на $3{,}18$\%, снижение коэффициента трения на $3{,}80$\% и рост расхода на $0{,}42$\%, что подтверждает положительное влияние параболического текстурирования на трибологические характеристики подшипника.
